\section{All overlaps from an arithmetic endpoint family}
\label{sec:overlaps}

\subsection{Candidate lengths forced by a short endpoint word}

For finite words $X,Y$, define
\[
 \mathcal O(X,Y)=\{k:1\leq k\leq\min(|X|,|Y|),\;
                    \suff_k(X)=\pref_k(Y)\}.
\]
These are literal signed-word equalities. Free reduction is not performed
inside this definition. In the application, the surrounding group algorithm
chooses the words whose overlaps are to be queried.

Write $L=\min(|X|,|Y|)$. Replacing $X$ by its last $L$ letters and $Y$ by
its first $L$ letters preserves $\mathcal O(X,Y)$. We therefore assume both
operands have length $L$. Pick $1\leq q\leq L$ and let
$P=\suff_q(X)$. Every overlap of length $k\geq q$ forces an occurrence of
$P$ in $Y$ ending at position $k$, where positions are counted from one.
Thus the set
\[
 K_P(Y)=\{i+q:Y[i:i+q]=P\}
\]
contains every possible overlap length at least $q$. This is only a necessary
condition. The algorithm must verify all lengths that it accepts.

An arithmetic progression is represented by $(a,d,c)$, denoting
$\{a,a+d,\ldots,a+(c-1)d\}$. For $c>1$, require $d>0$;
for a singleton, use $d=0$. Both $d$ and $c$ are exact binary integers.

\begin{theorem}[Endpoint progression clipping]
\label{thm:endpoint}
Suppose $K_P(Y)$ is a nonempty arithmetic progression $K$ with largest
member $M$, and put $U=\pref_M(Y)$. If $K$ has at least two members,
let its step be $d$ and suppose that $U$ has period $d$. Then
\begin{equation}
 \mathcal O(X,Y)\cap[q,L]
   =K\cap[q,\lcsuf(X,U)].
 \label{eq:endpoint}
\end{equation}
The hypothesis about the period is decidable by the exact equality
\begin{equation}
 U[0:M-d]=U[d:M].\label{eq:period-check}
\end{equation}
No periodicity hypothesis on the other operand $X$ is required.
\end{theorem}
\begin{proof}
If $K=\{M\}$, the only candidate is accepted exactly when
$\suff_M(X)=U$, equivalently when $M\leq\lcsuf(X,U)$.
Now suppose $K$ has at least two members.
Each $k\in K$ has the form $M-jd$ for an integer $j\geq0$.
The period condition implies $U[i]=U[i+jd]$ whenever both sides are within
$U$, by repeatedly applying the period identity. Therefore
\[
 \pref_k(Y)=\pref_k(U)=\suff_k(U).
\]
The condition $\suff_k(X)=\pref_k(Y)$ is consequently equivalent to
$\suff_k(X)=\suff_k(U)$, which holds precisely when
$k\leq\lcsuf(X,U)$. Conversely, every overlap of length at least $q$
belongs to $K$ by its final $q$ letters. These two implications prove \eqref{eq:endpoint}.
Equation \eqref{eq:period-check} is the definition of period $d$ for a finite word, so it gives
an exact certificate for the only periodicity used in the proof.
\end{proof}

\begin{corollary}[Symmetric endpoint direction]
\label{cor:endpoint-symmetric}
Let $Q=\pref_q(Y)$ and let
$K=\{L-i:X[i:i+q]=Q\}$. Suppose $K$ is a nonsingleton arithmetic
progression of step $d$ and maximum $M$, and $V=\suff_M(X)$ has period $d$.
Then
\[
 \mathcal O(X,Y)\cap[q,L]=K\cap[q,\lcp(V,Y)].
\]
\end{corollary}
\begin{proof}
For $k=M-jd$, periodicity gives
$\suff_k(X)=\suff_k(V)=\pref_k(V)$. Now use the definition of a longest
common prefix. The occurrence of $Q$ at the beginning of any overlap proves
necessity of membership in $K$.
\end{proof}

Empty and singleton candidate sets need separate, simpler treatment. If the
set is empty, no overlap of length at least $q$ exists. For a singleton,
one exact suffix--prefix comparison decides its membership. Every length
below $q$ is checked directly in all three cases. Hence one successful
anchor returns at most $q$ disjoint progression records: at most $q-1$
short singletons and one progression for the longer overlaps.

\subsection{Finding the cutoff by candidate rank}

The common-suffix and common-prefix lengths in the preceding statements
describe the answer, but need not be computed. A binary search for the
complete common-prefix length can examine positions that are irrelevant
to the progression. Once the progression and its period have been checked,
we can search only its candidate ranks.

\begin{corollary}[Galloping search for the last successful candidate]
\label{cor:endpoint-rank}
Under the hypotheses of \cref{thm:endpoint} or
\cref{cor:endpoint-symmetric}, write
$K=\{a+jd:0\leq j<c\}$ with $c\geq2$. Let $h$ be the number of its
members that are actual overlaps and let $z=c-h$.
All accepted members can be returned using
\begin{equation}
 O\bigl(1+\log(1+z)\bigr)
 \quad\text{exact suffix--prefix equality comparisons}.
 \label{eq:rank-comparisons}
\end{equation}
This is $O(\log(c+1))$ in the worst case. If every candidate succeeds,
one comparison suffices; if every candidate fails, at most two suffice.
These counts exclude the preceding period check and the fewer than $q$
short-length checks.
\end{corollary}
\begin{proof}
By the clipping theorem, the predicate
\[
 f(j)=\bigl[\suff_{a+jd}(X)=\pref_{a+jd}(Y)\bigr]
\]
is a true prefix: it holds precisely for $0\leq j<h$.
Check $f(c-1)$ first. If it holds, return all of $K$.
Otherwise check $f(0)$; if that fails, return the empty progression.
The remaining case has $1\leq h<c$. Maintain a checked true index
$\mathrm{lo}$ and a checked false index $\mathrm{hi}$, initially
$0$ and $c-1$.

Probe backward from the maximum at distances
$\delta=1,2,4,\ldots$, using $j=\max(0,c-1-\delta)$.
A false answer updates $\mathrm{hi}=j$; a true answer updates
$\mathrm{lo}=j$ and ends this phase. If $j=0$, its previously checked
true value supplies the endpoint without another comparison.
For an unclamped probe,
\[
 f(c-1-\delta)\text{ holds}
 \quad\Longleftrightarrow\quad \delta\geq z.
\]
Consequently the gallop uses $O(1+\log z)$ comparisons. Its final
bracket has width $O(z)$, including the clamped case. Binary search within
that bracket finds the last true index using $O(\log(1+z))$ additional
comparisons. The answer is the first $h$ members of $K$, represented by
one progression, or by a singleton when $h=1$.
The invariant also covers $c=2$, when no interior search is needed.
\end{proof}

The monotonicity is a consequence of the certified period condition.
Arbitrary overlap lengths need not have this property. The implementation
therefore applies the rank search only after the exact structural check;
it does not replace a failed period test by an assumed monotonicity.
The proof uses ordinary exponential search followed by bisection, rather
than a new general-purpose search method. Its useful feature here is the
choice of the candidate rank and the direction of the gallop.

\begin{example}[The opposite word may be nonperiodic]
Let $B$ be a word whose final letter $c$ appears nowhere else in $B$.
For integers $C\geq r\geq1$, let $Y=B^C$ and let $X=WB^r$, where
$W$ is any word of length $(C-r)|B|$. The endpoint occurrences in
$Y$ lie at $|B|,2|B|,\ldots,C|B|$. The theorem returns exactly those
multiples no larger than the actual common-suffix length. It does not need
to certify that $W$ repeats $B$, and it remains correct when $W$ contains
arbitrary defects or a partially matching suffix.
\end{example}

\subsection{Recognizing a progression without listing its points}

The missing operation is to summarize all occurrences of an explicit short
pattern in an SLP word. It is insufficient to summarize distinct grammar
nodes containing the pattern: a shared node can be used exponentially many
times in the represented word.

For an ordered finite set of occurrence starts $S$, retain
\[
 \Sigma(S)=(c,f,l,\delta_{\min},\delta_{\max}),
\]
where $c=|S|$, $f,l$ are its endpoints, and the last two entries are the
minimum and maximum consecutive gaps. Empty and singleton sets have both
gap entries zero. For $c\geq2$,
\begin{equation}
 S\text{ is an arithmetic progression}
 \quad\Longleftrightarrow\quad
 \delta_{\min}=\delta_{\max}.\label{eq:ap-test}
\end{equation}

\begin{lemma}[Exact ordered concatenation]
\label{lem:summary}
Suppose every point of $S$ precedes every point of $T$. The summary of
$S\cup T$ is determined by their summaries: counts add, the first and last
nonempty endpoints persist, and the new gap extrema are the extrema among
the old gaps and $\min T-\max S$. Translating $T$ only changes its two
endpoints. These formulas remain exact when the sets are exponentially large.
\end{lemma}
\begin{proof}
All consecutive pairs in $S\cup T$ lie within $S$, within $T$, or are the
single pair bridging their endpoints. This exhausts the gap list.
\end{proof}

For a rule $A\to BC$ and pattern length $q$, every occurrence is in exactly
one of three disjoint sets: wholly inside $B$; crossing the concatenation
cut; or wholly inside $C$. The last set is translated by $|B|$.
This locality at a grammar concatenation is standard in compressed
$q$-gram algorithms; see Goto, Bannai, Inenaga and Takeda
\cite{goto-qgrams}. The five-entry summary used here is tailored to
certifying one exact occurrence progression.
There are at most $q-1$ crossing positions. Store the first and last
$\min(q-1,|A|)$ letters of every reachable rule, and check those crossing
positions explicitly. Processing them in increasing start order allows
\cref{lem:summary} to merge all three sets without enumerating any child
occurrence set. For $q=1$ there are no crossing occurrences.

\begin{theorem}[SLP occurrence summaries]
\label{thm:summary}
For an explicit nonempty pattern of length $q$ and an SLP with $g$ reachable
rules, the occurrence summary can be computed using
$O(qg)$ stored letters and integers and
$O(q^2g+g\log(g+1))$ elementary operations with the implementation's
topological sorting. If $b$ bounds the bit length of positions, counts,
rule identifiers and signed letters, a conservative corresponding bit bound
is $O((q^2g+g\log(g+1))b)$, besides the input grammar's stored structure.
\end{theorem}
\begin{proof}
Terminal rules are immediate. The three-way partition at a concatenation
and \cref{lem:summary} establish correctness by induction on the grammar.
The retained endpoints have length at most $q-1$. There are at most $q-1$
crossing candidates, each checked against $q$ explicit letters, yielding
$O(q^2)$ elementary work per rule. The count for a repeated child is
included once on each side of the cut, correctly accounting for grammar
sharing. All integer additions, subtractions, comparisons and translations
involve $O(b)$ bits. Sorting the reachable rule identifiers adds the displayed
term; node order already respects the grammar's acyclicity.
\end{proof}

The summary is useful even when \cref{eq:ap-test} fails: it certifies that this particular
anchor does not give a single progression. It is not a complete compressed
representation of a nonprogression set. The matcher therefore tries another
anchor or falls back, rather than discarding unrepresented candidates.

\subsection{Integration and local complexity}

The implementation tries anchor sizes $1,2,4,8$, in both directions. Its
public primitive accepts other finite lists of positive sizes. A proposed
size is clipped to the available word length. The integrated hook is reached
only for a query window of length at least 128, after the earlier uniform,
short-period and sparse-endpoint shortcuts. Thus successful earlier paths
retain their previous work. The new hook is deterministic and preserves
the old complete dyadic occurrence-table fallback.

For one nonsingleton successful anchor, the implementation performs one
period equality and the candidate-rank comparisons in
\cref{cor:endpoint-rank}. It does not compute the complete longest common
prefix or suffix. Let $g$ count the rules reachable from the input operands,
let $H$ bound their height, and let $b$ bound the bit lengths of all
positions, counts, rule identifiers and signed letters involved in the
query. Write $\mathsf E(G,b)$ for a bound on one exact equality call on
an augmented grammar with at most $G$ rules, including its assertion and
cache bookkeeping. To account also for persistent arena data, let $g_0$
bound the total number of grammar rules and pre-existing equality-cache
records at entry; in particular $g_0\geq g$. Define
\[
 J=1+\left\lceil\log_2(1+z)\right\rceil.
\]
Each candidate comparison slices the same original operands. A substring
adds $O(H)$ rules and does not increase their height, so the initial
cropping, endpoint extraction, period slices, candidate slices and short
checks together create at most $O(H(q+J))$ additional rules. The equality
implementation itself uses assertions on existing rules. We may therefore
take
\[
 G=g_0+O\bigl(H(q+J)\bigr).
\]
This deliberately counts persistent data that a reachable-rule summary
does not inspect but an equality-cache cleanup might visit.
Charging $O(b^2)$ for elementary binary multiplication in the calculation
$a+jd$, a conservative bound for the new promised-class route is
\begin{equation}
 O\!\left((q^2g+g\log(g+1))b
 +(q+J)\bigl[\mathsf E(G,b)+Hb+b^2\bigr]\right).
 \label{eq:overlap-cost}
\end{equation}
The $q$ term includes the fewer than $q$ directly checked short lengths.
For the fixed list $1,2,4,8$, failed earlier anchors contribute only a
constant number of summary passes, endpoint slices and period comparisons;
the same bound applies with the largest tried $q$ and a changed absolute
constant. All large numerical dependence is through encoded lengths and
the exact comparison cost. In particular, \cref{eq:rank-comparisons}
improves the \emph{number} of comparisons; it does not make each equality
independent of the represented words. No asymptotic bound is transferred
from a different compressed-string implementation.

All calls use the same arena work allowance and cancellation callback.
Successful complete answers are cached under the caller's original root
pair, even when the actual comparison was cropped. Candidate-rank filtering
uses direct suffix--prefix comparisons in both anchor directions.
An interruption raises the existing resource exception, with no partial
overlap set entered into the answer cache.

An unsuccessful speculative certificate has a bounded number of summary
passes for the fixed policy, but its exact equality query may still be
expensive. This is why the experiments include failed period tests and
phase-shifted long-period queries. We claim the promised-class improvement
expressed in \cref{eq:overlap-cost}, not uniform dominance over every old query.

\subsection{An explicit family beyond both old numerical caps}

\begin{proposition}[Unbounded period and marker multiplicity]
\label{prop:marker-family}
For distinct signed letters $a,b,c$, positive integers $N,C$, and
$W=((ab)^Nc)^C$, let $d=2N+1$. Then
\[
 \mathcal O(W,W)=\{d,2d,\ldots,Cd\}.
\]
The $q=1$ endpoint method returns this set as a single progression, using
grammar operations polynomial in $\log(N+1)+\log(C+1)$ for its constructed
SLP. Neither $d$ nor $C$ has a fixed numerical upper bound.
\end{proposition}
\begin{proof}
The final letter is $c$, whose occurrences are exactly at multiples of $d$.
Every overlap must end at such a position in the second copy. Conversely,
each such length is a whole number of identical blocks and is an overlap.
The longest candidate window is all of $W$ and has period $d$.
Repeated squaring gives a grammar with
$O(\log(N+1)+\log(C+1))$ rules, and the exact compressed comparisons are
polynomial in that encoded grammar size.
\end{proof}

Two binary-alphabet variants demonstrate why short words are useful anchors:
\[
 W_2=((ba)^Na)^C,\qquad W_4=((ab)^Naa)^C.
\]
The terminal word $aa$ identifies exactly the block ends of $W_2$.
The terminal word $abaa$ does so for $W_4$. For $N\geq2$ the resulting
overlap sets are, respectively, the multiples of $2N+1$, and
$\{1\}$ together with the multiples of $2N+2$. The latter isolated overlap
is exactly why lengths shorter than the anchor must be retained.

With $N=C=2^{500}$, a period itself needs about 502 bits and the number of
overlaps needs 501 bits. Explicitly listing the overlaps is impossible within
reasonable resources, but their three-integer progression representation is
small. This family is a controlled compressed-string benchmark. No theorem
here says that arbitrary hard unknot instances must generate this family,
or that every important relator overlap satisfies its hypotheses.
